\documentclass{article}
\usepackage[T1]{fontenc}
\usepackage{lmodern}
\usepackage{graphicx}
\usepackage{amsmath,amssymb}
\usepackage[a4paper,margin=2cm]{geometry}
\usepackage[absolute,overlay]{textpos}
\setlength{\TPHorizModule}{1mm}
\setlength{\TPVertModule}{1mm}
\usepackage[indonesian]{babel}
\usepackage{hyperref}
\setlength{\emergencystretch}{2em}
% unit: Halaman 2

\begin{document}

% === Halaman 2 (python/native) ===

2

•

Monoalpabhetic cipher: kelompok cipher substitusi yang mengganti satu huruf 

dengan satu huruf yang lain 

•

Caesar cipher adalah salah satu cipher abjad-tunggal dengan melakukan substitusi 

huruf berdasarkan hasil  pergeseran huruf-huruf  alfabet sejauh k huruf. Namun 

Caesar Cipher bukan satu-satunya cipher abjad-tunggal.

•

Secara umum,  kita dapat membentuk tabel subtitusi sembarang. Jumlah 

kemungkinan tabel substitusi yang dapat dibuat pada sembarang cipher abjad-

tunggal adalah sebanyak 


   26! = 403.291.461.126.605.635.584.000.000 

      karena ada 26! cara mempermutasikan 26 huruf alfabet.

A. Cipher Abjad-Tunggal (monoalphabetic cipher)

\end{document}
